\printedpage{135}
% head: A. Grothendieck : Dérivateurs, Ch. XII, Édition des Universités de Montpellier II et Paris VII

\begin{quote}
takes $W_X$ into $W_Y$, where $W_X \subseteq \Fl(\cM/X)$ \texttt{[is]} defined by

$u \in W_X \iff \sigma(U) \in W$ \ ($\sigma : \cM/X \longrightarrow \cM$ the ``source'' functor, or ``forgetting $X$''). % [sic] \sigma(U) -> \sigma(u), the arrow is called u
\end{quote}

(\footnote[120]{\textbf{NB} Give here the proposition: $\mathrm{Fib}_W$ \textit{stable under composition and under base change, and contains the isomorphisms, stable under direct factors.}}). Let us put, for every $S$ in $\cM$,
\[
\Ho^!_W(S) = (\cM/S)W_S^{-1},
\]
which plays the role of a category of relative $W$-homotopy types \textit{locally constant} over $S$. We would like to describe a contravariant law in $S$ for $\Ho^!_W(S)$, i.e. for
\[
f : S' \longrightarrow S
\]
an arbitrary arrow in $\cM$, define
\[
f^*_W : \Ho^!_W(S) \xrightarrow{?} \Ho^!_W(S'). \tag*{$(*)$}
\]
The difficulty is, of course, that the base change functor
\[
f^* : \cM/S \longrightarrow \cM/S'
\]
is \textit{not} in general compatible with the \msp{page 209} localizers $W_S$, $W_{S'}$. The notion of $W$-fibration thus serves in the first place to bring out the most obvious condition on an arrow $f : S' \longrightarrow S$ in $\cM$, for the base change functor to be compatible with the localizers $W_S$, $W_{S'}$ (\textit{that} is the definition), hence induces $(*)$ on the localized categories.

The \textit{``problem of homotopy fibers''} is precisely to define $(*)$ for \textit{every} arrow in $\cM$, not only for $W$-fibrations, in such a way that the $\Ho^!_W(S)$ form, for $S$ % [sic] tellle -> telle, letter doubled
variable, the fibers of a remarkable fibered category over $\cM$. When $\cM$ admits a final object, so that we can regard the arrows
\[
p : e \longrightarrow S \qquad (p \text{ as ``point''})
\]
in $\cM$ as ``points'' of $S$, the ``$W$-homotopy fiber'' of an arbitrary object $X/S$ of $\cM/S$ at the said point will by definition be the element of
\[
\Ho_W = \Ho_W(\cM) \simeq \Ho^!_W(e)
\]
defined by
\[
p^*(\ho^!_{W/S}(X/S)),
\]
\printedpage{136}
% head: A. Grothendieck : Dérivateurs, Ch. XII, Édition des Universités de Montpellier II et Paris VII
\msp{page 210}
where
\[
\ho^!_{W/S} : \cM/S \longrightarrow \Ho^!_W(S) \overset{\text{def}}{=} (\cM/S)(W_S)^{-1}
\]
is the canonical localization functor.

For every $S$ in $\cM$, let
\[
\Fib_W S \subseteq \cM/S
\]
be the full subcategory formed by the objects $X/S$ of $\cM/S$ such that $X \longrightarrow S$ is $W$-fibrant. We are going to define a variant
\[
\Ho^{\lc}_W(S)
\]
of $\Ho^!_W(S)$, as
\[
\Ho^{\lc}_W(S) \overset{\text{def}}{=} (\Fib_W S)(\Wu_S)^{-1},
\]
where

\begin{quote}
$\Wu_S \subseteq \Fl(\Fib_W S)$ is the localizer formed by the arrows $u : X \longrightarrow Y$ which are ``universally over $S$ in $W$'', i.e. such that for every arrow $S' \longrightarrow S$ of $\cM$ with target $S$, the arrow deduced by base change,
\[
u' : X' = X \times_S S' \longrightarrow Y' = Y \times_S S'
\]
is in $W_{S'}$.
\end{quote}

We have ``done what was needed'' so that $\Ho^{\lc}_W(S)$ is contravariant in $S$ for \textit{arbitrary} arrows $f : S' \longrightarrow S$ of $\cM$, they define
\[
f^{*\,\lc}_W : \Ho^{\lc}_W(S) \longrightarrow \Ho^{\lc}_W(S'), \tag*{$(**)$}
\]
\msp{page 211}
since the $W$-fibrations are stable under base change, and since consequently $f^* : \cM/S \longrightarrow \cM/S'$ induces
\[
f^*_{\Fib_W} : \Fib_W S \longrightarrow \Fib_W S'
\]
which, by definition of $\Wu_S$ and $\Wu_{S'}$, is compatible with the localizers
\[
f^*_{\Fib_W} : (\Fib_W S, \Wu_S) \longrightarrow (\Fib_W S', \Wu_{S'}),
\]
and thus passes to the categories of fractions as an arrow $(*)$.

We may wonder why not be more generous, and define $\Ho^{\lc}_W(S)$ as $(\cM/S)(\Wu_S)^{-1}$, defining $\Wu_S$ in $\cM/S$ in the same way as just now in $\Fib_W S$ - which costs no more. A first answer is from experience the localizer $\Wu_S$ in the whole of $\cM/S$ turns out to be practically unusable, on the one hand because membership in $\Wu_S$ (for an arrow between objects of $\cM/S$ which are not very special, e.g. $W$-fibered over $S$) is unverifiable (except for isomorphisms and things like that \dots), and because on the other hand the category of fractions is prohibitively
\printedpage{137}
% head: A. Grothendieck : Dérivateurs, Ch. XII, Édition des Universités de Montpellier II et Paris VII
\msp{page 212}
big and can serve (it seems) no purpose. Another answer is that we want the category obtained to correspond well to the intuition of relative homotopy types \textit{locally constant} over $S$ - and we therefore want that for an arbitrary base change $f : S' \longrightarrow S$ by an

\begin{quote}
$S'$ in $\cM/S$, whence a composite functor
\[
\begin{tikzcd}[column sep=small]
\cF_S(\Wu_S)^{-1} \arrow[r, "f^{*\,\lc}_W"] \arrow[rr, dotted, bend right=15, "\varphi_{S'}"'] & \cF_{S'}(\Wu_{S'})^{-1} \arrow[r] & \Ho_W
\end{tikzcd}
\]

(\footnote[121]{$\cF$ here denotes a fibered subcategory of $\cM/\cM = \underline{\Fl}(\cM)$ over $\cM$, a possible substitute for the one formed by $\Fib_W S$ or by the whole of $\cM/S$.}) deduced from the functors of categories with localizers
\[
(\cF_S, \Wu_S) \xrightarrow{f^*} (\cF_{S'}, \Wu_{S'}) \xrightarrow{\text{forgetting } S'} (\cM, W)
\]
\end{quote}

(composite functor $\varphi_{S'}$ which can be viewed as the functor \textit{``homotopy fiber at the object $S'$ over $S$''}, this object being visualized as a ``point'' of $S$ ``with values in $S'$ '' following the general functorial yoga \dots), the functor
\[
\begin{aligned}
S' &\longmapsto \varphi_{S'}(\xi) \qquad \text{for } \xi \text{ in } \cF_S(\Wu_S)^{-1}\\
\cM/S &\longrightarrow \Ho_W
\end{aligned}
\]
take arrows in $W_S$ to isomorphisms: the ``$W$-homotopy fiber'' at $S'/S$ of an object of $\cF_S(\Wu_S)^{-1}$ (viewed as a $W$-homotopy type \textit{locally constant} over $S$) must be ``the same'', if $S'$ replaced by $S''$ with $S'' \longrightarrow S'$ given and in $W_S$. It is for this reason that it is reasonable to take $\cF_S$ ``small enough'' (namely $\Fib_W S$, which is \msp{page 213} exactly adjusted for this), rather than take the whole of $\cM/S$.

We can therefore say that the category $\Fib_W S$, and the localizer $\Wu_S$ have been constructed ad-hoc, as simply as possible, to give rise to a localized $\Ho^{\lc}_W(S)$ which corresponds to the intuition of locally constant relative $W$-homotopy types in the following two ways:

\begin{enumerate}[label={}]
\item[1°)] There is spontaneously a contravariant law in $S$, for arbitrary arrows $S' \longrightarrow S$ in Cat (we have defined $\Wu_S$ just for this purpose accordingly).
\item[2°)] that the notion of ``$W$-homotopy fiber at $S'/S$'' of an object $\xi$ of $\Ho^{\lc}_W$ satisfy the preceding condition, i.e. an arrow $S'' \longrightarrow S'$ in $\cM/S$ which is in $W_S$ define an \textit{isomorphism} (in $\Ho_W$) of homotopy fibers
\[
\underbrace{\xi_{S''}}_{\varphi_{S''}(\xi)} \longrightarrow \underbrace{\xi_{S'}}_{\varphi_{S'}(\xi)}
\]
(and we have chosen $\cF_S = \Fib_W(S)$ just for that purpose).
\end{enumerate}

\printedpage{138}
% head: A. Grothendieck : Dérivateurs, Ch. XII, Édition des Universités de Montpellier II et Paris VII

We thus have an excellent notion of $W$-homotopy fibers for $X$ $W$-fibered over $S$ - this time excellent for \textit{arbitrary} base changes - and there is no merit in that since we have defined everything on purpose for that! % [sic] quelconque -> quelconques, agreement with "des changements"

\msp{page 214}
To put it otherwise, we have defined \textit{two} notions of relative $W$-homotopy types over $S$, one so as to give such a homotopy type $\ho^!_{W,S}(X)$ for \textit{every} object of $\cM/S$, but there is a priori a good notion of base change only for the base changes $S' \longrightarrow S$ which are $W$-fibered, the other, which corresponds more directly to the intuition of a locally constant $W$-homotopy type over $S$, associates such a type $\ho^{\lc}_{W,S}(X)$ only to the $X$ of $\cM/S$ which are $W$-fibered over $S$, on the other hand this notion admits a good notion of inverse image and of ``$W$-homotopy fiber'' for arbitrary base changes.

A ``good'' theory of $W$-homotopy fibers must combine the two preceding notions into a single one, so that the advantages of the two add up, and the disadvantages cancel out! It is obtained with the help of a draconian axiom on the situation (whereas until now we have used nothing about $W \subseteq \Fl(\cM)$, except at most weak saturation, in addition to the existence of fiber products):

\msp{page 215}
\circled{a} (\textit{Homotopy fiber axiom.}) For every object $S$ in $\cM$, the canonical functor
\[
\Ho^{\lc}_W(S) \longrightarrow \Ho^!_W(S)
\]
deduced from the inclusion
\[
(\Fib_W(S), \Wu_S) \longrightarrow (\cM/S, W_S)
\]
is an \textit{equivalence of categories.}

This axiom implies that in $\Fib_W S$, $\Wu_S$ and the localizer induced by $W_S$ (which I still denote by $W_S$ by abuse of notation) have the same strong saturation. In fact, it is natural to supplement the preceding axiom with the following one.

\circled{b} (Adjunct axiom to the homotopy fiber axiom.) On $\Fib_W S$, the localizers $\Wu_S$ and $W_S$ coincide, i.e. we have this:

\begin{quote}
Let $X$, $Y$ be in $\Fib_W S$, $X \xrightarrow{u} Y$ an $S$-morphism such that $u \in W_S$, then $u \in \Wu_S$ (i.e. the morphism $u' : X' \longrightarrow Y'$ deduced by any base change $S' \longrightarrow S$ is still in $W$).
\end{quote}

\msp{page 216}
I presume that a) will rarely be satisfied without b) being so (and conversely for that matter), and that the proof of a) and b) will be done practically in a single breath.

If $\cM = \mathrm{Cat}$, $W$ a basic localizer, then the notion of $W$-fibration adapted here (read: ``ultrastrong'', in the specific terminology developed in Cat) is so strong, that it already practically implies b) (assuming W(4a$'$), it is true, cf. p. 53), and
\printedpage{139}
% head: A. Grothendieck : Dérivateurs, Ch. XII, Édition des Universités de Montpellier II et Paris VII
a) is hardly verifiable unless we have W(4a) (\footnote[122]{No, it seems to me that W(4b) should suffice \dots}) (loc. cit.). So here it is ``the fibration axiom'' on $W$, in its strong form W(4a), which is the keystone of a good theory of homotopy fibers.

Such a theory (returning to the general case) amounts to saying, in essence, that there are ``enough'' $W$-fibrations (``ultrastrong'') to be able to describe \textit{every} object of a $\cM/S$, up to $W_S$-equivalence, as a $W$-fibrant object over $S$. In the cases known to me, we have moreover this: we \msp{page 217} have for every $S$ in $\cM$ a subcategory $\cF_{0\,S}$ of $\Fib_W S$ (I am thinking for example of $\underline{\mathrm{Parf}}_W S$, or even $\underline{\mathrm{Parf}}_{W_\omega} S$, in the case where $\cM = \mathrm{Cat}$), such that for every full subcategory $\cF$ of $\cM/S$ such that
\[
\cF_{0\,S} \subseteq \cF \subseteq \cM/S
\]
(a fortiori, for a full $\cF$ such that $\Fib_W S \subseteq \cF \subseteq \cM/S$), the inclusions induce equivalences of categories
\[
\cF_{0\,S}(W_S)^{-1} \xrightarrow{\sim} \cF(W_S)^{-1} \xrightarrow{\sim} \underbrace{(\cM/S)(W_S)^{-1}}_{\overset{\text{def}}{=} \mathrm{Hot}^!_W(S)}. % [sic] Hot -> Ho, elsewhere the notation is Ho^!_W(S)
\]
It does not seem possible to fix the choice of $\cF_{0\,S}$ in ``the best possible'' way, once and for all in the axiomatic context adopted here. It is an additional element of structure in $\cM$ (besides the localizer $W$), on the same footing as the datum of the classes of $Q$-fibrations and $Q$-cofibrations, in the manner of \textsc{Quillen}. I do not know whether the axioms a), b) above (homotopy fiber axiom) already imply \msp{page 218} the variant of \circled{a}:

\circled{c} For every full subcategory $\cF$ of $\cM/S$ such that
\[
\Fib_W S \subseteq \cF \subseteq \cM/S
\]
the inclusions induce equivalences
\[
\underbrace{(\Fib_W S)(W_S)^{-1}}_{\begin{subarray}{l} = \ \Ho^{\lc}_W(S)\\ \text{assuming a)}\end{subarray}} \xrightarrow{\sim} \cF(W_S)^{-1} \xrightarrow{\sim} \underbrace{(\cM/S)(W_S)^{-1}}_{=\Ho^!_W(S)}.
\]

But I presume that in practice, when we can verify \circled{a}, we will be able to verify \circled{c} as well (\footnote[123]{Cf. page 240 ff. We should here develop a little more the formalism of $\Ho^{\lc}_W(S)$ and homotopy fibers, under the hypothesis \circled{a}\,\circled{b} \ above.}).